\section{Synthetic Evaluation}

The first RCOUNT milestone uses synthetic fixtures instead of real election
exports. This is a conscious staging decision. Real data adapters introduce
jurisdiction and vendor variation before the core contract is stable. Synthetic
fixtures let each verifier layer be exercised with minimal records, known
expected outcomes, and reproducible failure modes.

\begin{center}
\small
\begin{tabularx}{\textwidth}{llX}
\toprule
Fixture & Expected & Contract layer \\
\midrule
\texttt{summary-basic} & pass & Contest sums and jurisdiction total \\
\texttt{canvass-correction} & pass & Status transition and correction event \\
\texttt{mail-batch-added} & pass & Batch accounting and late mail evidence \\
\texttt{precinct-split-lineage} & pass & Split and merge reporting-unit lineage \\
\texttt{privacy-inclusion-sketch} & pass & Choice-free inclusion proof sketch \\
\texttt{district-aggregation-rplan} & pass & Optional RPLAN district aggregation \\
\texttt{multi-election-harness} & pass & Three-cycle L2 replay with lineage and districts \\
\texttt{bad-selection-sum} & fail & Local contest arithmetic \\
\texttt{missing-batch} & fail & Missing batch evidence \\
\texttt{bad-lineage} & fail & Missing current lineage unit \\
\texttt{choice-bearing-proof} & fail & Receipt-risk proof exposure \\
\texttt{tampered-source} & fail & Source hash mismatch \\
\texttt{missing-source-hash} & fail & Omitted source evidence \\
\texttt{multi-election-harness-negatives} & fail & Bad lineage, stale plan, tampered cycle source \\
\bottomrule
\end{tabularx}
\end{center}

\begin{center}
\small
\begin{tabularx}{\textwidth}{llX}
\toprule
Level & Scope & Evidence in this draft \\
\midrule
L0 & Pure equations and hash projections & Core unit tests for contest sums, jurisdiction totals, status events, package hashes, batch accounting, lineage, and proof privacy. \\
L1 & Package fixtures and CLI verification & Golden packages under \texttt{docs/examples/rcount-golden-packages} and \texttt{rcount verify}. \\
L2 & Multi-election replay & Three-cycle synthetic harness with split/merge lineage, per-cycle RPLAN plans, district aggregation, and negative cases. \\
\bottomrule
\end{tabularx}
\end{center}

The L2 harness is the most important fixture because it links the concepts that
otherwise look separate. It generates three synthetic election cycles. The 2026
cycle splits one precinct into two. The 2028 cycle merges two precincts into a
new reporting unit. Each cycle has a package, a plan assignment, and district
aggregation evidence. Negative variants then break the lineage, use a stale
RPLAN unit, or tamper with one cycle's source bytes.

This style of synthetic evaluation is not a claim about performance or
real-world coverage. It is a contract test. Before RCOUNT reads a public state
file, the package format should already know how to fail when arithmetic,
lineage, privacy, district projection, or source evidence breaks.

The reproducibility surface is executable. The key commands are:
\begin{quote}
\small
\begin{verbatim}
cargo run -p rcount-district --example multi_election_harness
cargo run -p rcount-district --example multi_election_negative_harnesses
cargo test -p rcount-district -p rcount-audit -p rcount-cli
\end{verbatim}
\end{quote}
Negative fixtures intentionally fail through different surfaces:
\begin{itemize}
\item bad lineage fails the package verifier;
\item stale plan units fail district aggregation;
\item tampered cycle sources fail source hash verification.
\end{itemize}
